\documentclass{ctexart}
\usepackage{amsmath,amssymb}
\begin{document}
\section*{中断下含短折返的轨道交通恢复调度}
当前有限事件模型补齐原文式(16)--(22)的恢复终止和费用，保留
旧02的时刻、客流、发车间隔与车辆可用性约束。源文图9的完整16站、
15段线路和47辆物理列车均保留。未恢复的需求及初态采用已有补全数据。
当前代码通过策略搜索生成可行调度。

\subsection*{输入与决策}
两方向$r\in\{0,1\}$，中断$[T_0,T_0+\widetilde T)$。
控制节点在阻断区段两侧。$\eta_{rk}$为第$k$项服务控制节点出发时刻，
$h_{rk}$为其间隔，$H=\{100,150,\ldots,500\}$秒；
正常参考时刻表分别171、180班，定义恢复参照而非强制重排351班。
站间运行时间$v$、停站$\omega=25$秒、折返$\Delta=180$秒、
车库整备及容量$D$均由输入给定。
车辆具有唯一标识、当前位置及最早可用时刻。
决策是有限服务序列、各$h_{rk}$、合法路线和实际车辆分配。
\[
\eta_{rk}=\eta_{r,k-1}+h_{rk},\qquad h_{rk}\in H.
\]
前$K'_r$列为初始受阻返回服务；后续在
$\eta_{rk}-\Delta-\omega<T_0+\widetilde T$时选择同侧短折返环线，
否则为正常全线。沿路径相邻不同站
$\tau_{n',k}=\tau_{n,k}+v_{nn'}+\omega$；
折返点重复站增加$\Delta+\omega$。堵塞解除前不得进入阻断区段。

\subsection*{车辆与折返资源}
记输入初始最早折返时刻$\bar\tau_{rk}$，初始返回服务满足
\begin{align*}
\eta_{rk}&\ge\bar\tau_{rk}+\Delta+\omega,\\
\eta_{rk}&\ge\bar\tau_{r,k-1}+2\Delta+\omega &&k>1,\\
\eta_{rk}&\le\bar\tau_{r,K'_r}+H^{max}+\Delta+\omega
 -(K'_r-k+1)H^{min}.
\end{align*}
折返区间为$[\eta-\omega-\Delta,\eta-\omega]$，同侧不得重叠。
后续环线满足$h\ge\Delta$及实际始发时刻不早于车辆可用时刻。
正常恢复服务也满足车辆可用性和解堵条件。
每辆物理车连续两服务的前者终点必须等于后者始点，
且前者终到加整备不晚于后者始发；初始库存按输入登记，
可用车辆按先入先出使用。该显式车辆链代替旧02式(30)/(32)可能出现
负下标的全局服务编号，并按物理车辆集合核对。
同方向同站相邻服务出发至少间隔100秒。

\subsection*{完整OD队列与列车容量}
状态包含每方向、起终点的持续队列、车上OD人数、车辆位置/可用时刻、
已完成事件及正常参照。初始站台队列来自完整正常时刻表运行至$T_0$；
初始受阻车在边界疏散后空车返回，这是明确的补全初态。
到达率$\widetilde O_{r,o,d}(t)$按分钟分段常数给定；
中断期不可达新增跨阻断需求按源文S1转其他交通，之前积压仍保留。
每次站台服务事件$t$，上次同方向同站服务为$t^-$：
\[
q_{rod}(t)=l_{rod}(t^-)+\int_{t^-}^{t}\widetilde O_{rod}(u)\,du .
\]
先令车上终点为$o$的乘客下车，再定义本次路线后续可达的目的地集合$E$。
剩余容量为$D-B$，登乘按源文比例分配：
\[
\gamma=\begin{cases}
\min\{1,(D-B)/\sum_{d\in E}q_{rod}\},&\sum_{d\in E}q_{rod}>0,\\
0,&\text{否则},
\end{cases}\quad
\widehat q_{rod}=\mathbf1_{\{d\in E\}}\gamma q_{rod},\quad
l_{rod}=q_{rod}-\widehat q_{rod}.
\]
车上人数按目的地守恒更新，始终在$[0,D]$内；站台队列非负。
短折返只能运送同一连通侧且当前行驶方向可达的OD。

\subsection*{有限终止与目标}
每个方向结束于一班与正常参照具有相同完整轨迹的服务，
且该服务各站出发后剩余队列不高于对应正常参照，恢复条件对应原文(16)。
服务结束时车上乘客全部到达；未上车队列继续保留并报告。
目标为最大化有限恢复序列的负总费用：
\[
\max C=-\left[
\frac{c_W}{3600}\sum_{\text{站台事件 }e}\int_{t_e^-}^{t_e}
\sum_d Q_{r_e,o_e,d}(u)\,du
+\sum_{\text{行驶弧 }e}c_E\epsilon_e(m_T+m_PB_e)
+\sum_{\text{行驶弧 }e}(c_D L_e+c_T(\tau'_e-\tau_e))
+c_S K_{\rm short}\right].
\]
等待积分覆盖每个站台从初态/上次事件至本次事件的排队面积，
各方向的费用累积区间截止于已验证的恢复服务时刻。
运行能耗按弧的载客重量计算，转向重复站不计不存在的行驶距离，
每项初始返回/短折返服务计一次短折返成本。

\subsection*{算法与交付证据}
代码搜索12组偏好间隔，在$H$内按可用性调整，并在恢复阶段接入正常轨迹。
搜索方法使用12组偏好生成候选策略。
输出所有事件、OD队列、车辆链、费用明细与恢复证书，逐项复算容量、
持续队列、车辆唯一性、阻断、折返和恢复条件。
当前入口在函数/循环边界执行3600秒预算检查，包括尚未找到可行候选时；
结果记录所生成可行策略的目标值。
\end{document}
